% AI-produced: GPT-6 Astra (OpenAI), /root/axioms, 19 November 2026.
% Research contribution, with no priority claim. Parent manuscript supplies
% theorem environments, natbib and standard AMS packages.
\section{B.1. A typed comparison of the successive axiom systems}
\label{SA21-B-axioms}

We distinguish the written axioms from the stronger hypotheses tacitly used
in several of their purported consequences. In particular, complete
join-irreducibility is not complete join-primality. This distinction changes
the comparison already at C1--C3, before any refinement is introduced.
The symbols C4 and C8 also change their meanings during 156-4. We therefore
attach a scan locator whenever we use a chronological label.

\begin{definition}[Known: notation and the source-resolved prime convention]
\label{SA21-B-notation}
Let \(\mathscr F\) be the poset of figures and let \(M\) be its nonzero
completely join-irreducible elements. Write
\[
\begin{aligned}
 \widetilde F&=\{X\in M:X\leq F\},\\
 F\bowtie G&\Longleftrightarrow\{F,G\}\text{ has an upper bound},\\
 F\parallel G&\Longleftrightarrow F\bowtie G\text{ and }F\wedge G=0.
\end{aligned}
\]
We call the incidence system \emph{prime} if every \(X\in M\) is
completely join-prime for every bounded family:
\[
 X\leq\bigvee_{i\in I}F_i\quad\Longrightarrow\quad
 X\leq F_i\text{ for some }i\in I.                 \tag{B.P}
\]
After proving the representation below, and only in the prime case, we
identify \(F\) with its downward-closed deployment \(\widetilde F\).
Thus \(F\cap G\) is incidence intersection, not intersection of supports.

The refinement order is \(\ll\), with \(\leq\subseteq\ll\).
For \(X\ll F\), its \emph{carrier} is the least \(Y\in\widetilde F\)
such that \(X\ll Y\). We write \(X\rho Y\), or \(X\ll^\circ Y\), if
\(Y\) is its carrier in \(\widetilde Y\). Lieux are the
\(\ll\)-minimal members of \(M\), and
\[
 \ell(F)=\{p\in L:p\ll F\},\qquad
 U_X=\{p\in L:p\rho X\}=\ell(X)\setminus\ell(\partial X),
 \quad \partial X=\widetilde X\setminus\{X\}.
\]
For refinements of subfigures we use
\[
 P|G=\{Z\in\widetilde P:Z\ll G\},\qquad
 P\preceq F\Longleftrightarrow P\ll F\text{ and }
 (H\parallel P\Longleftrightarrow H\parallel F)\text{ for every }H.
                                                        \tag{B.S}
\]
Primitive internal disjointness on \(M\) is denoted \(\delta\), the
manuscript's \(\mid\circ\mid\). It must not be confused with figure
disjointness \(\parallel\).
\end{definition}

\subsection{The literal incidence axiom and its exact repair}

\begin{theorem}[Proved: exact bounded-join representation criterion]
\label{SA21-B-prime-representation}
Suppose every bounded family in the nonempty poset \(\mathscr F\) has a
join, and every figure is the join of its completely join-irreducible
strata. Then \(F\mapsto\widetilde F\) is an order embedding and preserves
all nonempty meets. The following are equivalent:
\begin{enumerate}
\item it preserves every bounded join as a union;
\item all strata satisfy \textup{(B.P)};
\item its image contains every downward-closed subset of every one of its
members, with joins there given by unions.
\end{enumerate}
Under these equivalent conditions, compatibility by strata is precisely
the condition that two admitted downsets have an admitted union whenever
all their cross-pairs of principal ideals have admitted unions. If every
figure has finitely many strata, the prime condition is equivalent to
distributivity of every principal figure ideal.
\end{theorem}

\begin{proof}
Generation gives \(F\leq G\) exactly when every stratum below \(F\) is
below \(G\); hence the order embedding. Every nonempty family has a meet:
take the bounded join of its common lower bounds. A stratum lies below
that meet exactly when it lies below each family member.
The equality
\(\widetilde{\bigvee_iF_i}=\bigcup_i\widetilde F_i\)
is precisely (B.P). If \(D\subseteq\widetilde F\) is a downset,
\(H=\bigvee_{X\in D}X\) exists; primality gives
\(\widetilde H=\bigcup_{X\in D}\widetilde X=D\).
Conversely, an image consisting of downsets and closed under bounded
unions has those unions as its joins, because the deployment map is an
order embedding. This proves the equivalences and the compatibility
translation.

For the final assertion, every interval \([0,F]\) is finite. A finite
downset lattice is distributive. Conversely, in a finite distributive
lattice a join-irreducible element \(X\) satisfying \(X\leq A\vee B\)
obeys \(X=(X\wedge A)\vee(X\wedge B)\); irreducibility puts it below one
of \(A,B\). Induction treats finite joins, and every family in a finite
lattice has the same join as a finite subfamily. An irreducible of
\([0,F]\) is also irreducible in \(\mathscr F\), since any expression
with join \(X\leq F\) has all its terms below \(F\). This proves (B.P).
\end{proof}

\begin{counterexample}[Refuted: literal C1--C3 and the literal atlas Scholie]
\label{SA21-B-M3}
Let \(\mathscr F=M_3=\{0,a,b,c,1\}\), with three incomparable atoms and
\(a\vee b=a\vee c=b\vee c=1\). Put \(\ll=\leq\).
All bounded joins exist; the completely join-irreducibles are
\(a,b,c\); every figure is their join; all pairs are compatible.
Thus literal C1--C3 hold. Refinement by strata holds because deployment
inclusion detects the order. The carrier of an atom in any figure
containing it is that atom; compatibility plus refinement implies
incidence; internal refinement on \(M\) is equality. Thus the displayed
At 1--At 7 clauses also hold when \(M\) is defined by the
equality-based irreducibility of 156-4 and 156-8. Nevertheless
\[
 \widetilde{a\vee b}=\{a,b,c\}\ne\{a,b\}
   =\widetilde a\cup\widetilde b.
\]
The downset \(\{a,b\}\) of \(\widetilde1\) is not a deployment.

This example survives the literal atlas axioms of 156-8, scans 20--21.
On \(M=\{a,b,c\}\), take incidence and refinement to be equality and
\(x\delta y\) to mean \(x\ne y\). Mag 1--4 hold: all principal
deployments are singletons, and internal heredity is trivial. Atlas
At 1 is bounded completeness, At 2 is generation, At 3 is total
compatibility, At 4 is the deployment test for refinement, and At 5(a)
is satisfied by defining \(\delta\) only on elementary objects.
At 5(b) holds because distinct atoms are \(\delta\)-related. Thus
\emph{even atlas At 1--5 plus Mag 1--4 does not imply the advertised
downset representation}. Adding (B.P) is necessary and sufficient.
\end{counterexample}

The 156-4 source defines irreducibility by an equality
\(F=\sup_iF_i\), not by \(F\leq\sup_iF_i\).
The unproved passage from joins to unions occurs immediately afterward
\citep[156-4, scans 54--57]{Grothendieck1986}:
\path{J156-4-07_transcription.tex:84--115,153--155}.
The repeated atlas claim is
\path{J156-8-03_transcription.tex:106--168}.
There is a crucial chronological qualification: formula (1.20) on
156-6/12 explicitly uses \(X\leq\sup_iF_i\Rightarrow
X\leq F_i\) for some \(i\), namely (B.P), before scan 17 renames
the old C1--C3 as At 1--At 3. We verified the inequality on the
original scan; see
\path{J156-6-02_transcription.tex:77--84,262--265}.
Thus this counterexample does not refute the full 156-6 presentation
when its explicit definition (1.20) is retained.
The 156-7 phrase ``strictement irr\'eductibles'' at
\path{H156-01_transcription.tex:64--67} gives no new formula:
the inherited 156-6 prime meaning validates the representation,
whereas the 156-4 equality meaning does not.
The finite certificate \texttt{literal\_systems.py} checks every subset
of \(M_3\), identifies its irreducibles, verifies the deployment tests,
and checks Mag 1--4 on the induced three-object structure.

\subsection{Chronological normalization}

All source filenames in this subsection are relative to
\path{RESEARCH_SA21_transcriptions/transcriptions/156_Codex_web/}.
We use the following source-resolved dictionary; the qualifications are
part of the mathematical assertions.

\begin{longtable}{p{.17\linewidth}p{.74\linewidth}}
\hline
Source and labels & Exact content and comparison\\ \hline
156-4/18--25, Contr 1--5 & Incidence is refinement plus compatibility;
finite geometric dimension; nonempty open lieu-shadows; disjoint open
shadows for distinct compatible strata; finally a proposed converse
recovering refinement and compatibility from open shadows. The last
converse is an additional set-theoretic reflection condition, already
qualified in the manuscript.\\
156-4/26--31, Cont 1--4 & Set-theoretic figures, with the empty figure and
point figures, the preliminary admissibility clauses, closure under closed subfamilies, and
compatibility tested on elementary figures. This is a different
numbering and is not an abstract finite-dimension axiom system.\\
156-4/34--44, Cont 1--5 & The incidence order and hereditary
compatibility; Fig 1--4 specify principal figures, clique downsets,
closed subfigures and compatible unions; Cont 3 is finite geometric
dimension, with optional Cont 3\( '\) finite principal deployments;
Cont 4 is nonempty open lieu-shadows and their setwise separation;
Cont 5 is disjointness hereditary under refinement. Optional Cont 6
detects disjointness on lieux; Cond 7 is their orthogonality across
distinct compatible open strata; Cond 8 is induced refinement.\\
156-4/53--83, C1--C8\(_0\) & C1 bounded joins; C2 irreducible generation;
C3 compatibility by strata; C4 compatibility across distinct incidence
components; C5 subdivisions preserve cosupport; C6 induced subdivisions,
with uniqueness, the stratum formula and join preservation;
C7\(_0\) extension of subdivisions; C8\(_0\) excludes an incidence and
subdivision cycle through an elementary object. C4 bis asks for
componentwise assembly. Final C8, at scan 87, is a separate gluing axiom.\\
156-6, At 1--7 & Formula (1.20) on scan 12 explicitly defines strata
by bounded-join primality, before scan 17 renames C1--C3.
With that definition the incidence presentation is prime; the later
clauses are refinement by strata, least carrier, compatible refinement
implies incidence, and internal transitivity. Importing only 156-4's
equality-based irreducibility gives a weaker presentation.\\
156-7, At 1--7 & Same intended core, but At 5 and At 6 exchange numbers.
The written hypothesis \(X\triangleleft F\) in At 6 must be
\(X\ll F\), as in its own \(M\)-translation. The term
``strictement irr\'eductibles'' inherits the prime meaning of 156-6
if that earlier definition is retained.\\
156-8, Mag 1--4 & Two orders with \(\leq\subseteq\ll\); least carriers
in principal deployments; symmetric irreflexive internal disjointness
hereditary under internal refinement; all distinct strata under one
top are internally disjoint.\\
156-8/30--32, primed and div & Internal refinement becomes primitive.
The strong primed version is equivalent to Mag 1--4. The final written
Mag 4\( '\), and later Mag div 3, require only comparable pairs and are
strictly weaker. Divisibility and the two-strata clause repair the
local formulation.\\
156-8/20--21, atlas At 1--5 & Bounded joins, generation, compatibility
and refinement detected on strata, and the primitive internal
disjointness formula, together with Mag 1--4. Add (B.P) for the claimed
equivalence with an actual atlas of clique downsets.\\
156-9/41, At 1--4 & The change of names does not change the strong
relational theory. A comparable-only reading inherits the earlier
weakening.\\
156-9/47--50, At spat 1--3 & An incidence order, internal disjointness,
the comparable-pair clause, and an injective deployment into
\emph{ordered spatial figures}. The target's pairwise orthogonality
condition is essential; with that typed target the strong two-strata
axiom is restored. Ordinary disjoint-set labels do not suffice.\\ \hline
\end{longtable}

For Cont, the exact successive lists are
\path{J156-4-03_transcription.tex}, scans 18--25;
\path{J156-4-04_transcription.tex}, scans 26--31; and
\path{J156-4-05_transcription.tex:6--28,43--48,71--109,202--211},
scans 34--40.
For the later core see
\path{J156-6-03_transcription.tex}, scans 19--25, and
\path{H156-01_transcription.tex:139--197}.
The printed At 6 slip is recorded in the input correction
SA21-K7-01; we use the explicit \(M\)-formulation at lines 188--197,
which removes the ambiguity. Read literally with \(X\triangleleft F\),
that clause follows from compatible refinement being incidence and
does not require a carrier for an external refinement. The finite
model in Counterexample~\ref{SA21-B-C-carriers} satisfies this literal
clause, At 4, compatible-refinement incidence, and internal transitivity,
but has no least carrier for \(a\ll dyz\). Its certificate checks
each of those additional assertions. Thus the corrected and literal
readings are genuinely different.

\begin{counterexample}[Refuted: the preliminary infinite admissibility shortcut]
\label{SA21-B-preliminary-admissibility}
The preliminary set-theoretic Cont 2 at 156-4/27 requires nonempty
residuals and binary-intersection coverage; for infinite families
these do not imply the later admissibility definition. On
\(L=\mathbb N\cup\{\infty\}\), put
\(A_n=\{\infty\}\cup\{m\in\mathbb N:m\geq n\}\).
The family \(\{A_n:n\in\mathbb N\}\) has
\(A_n^\circ=\{n\}\) and
\(A_n\cap A_m=A_{\max(n,m)}\), but \(\infty\) belongs to no
residual and has no least stratum. Adding coverage by the residuals
is necessary and sufficient: binary-intersection coverage makes
any minimal containing stratum least. The later text explicitly
adds that coverage condition
\citep[156-5, scans 2--5]{Grothendieck1986},
\path{J156-5-01_transcription.tex:22--37,151--154}.
All uses of \(\operatorname{Fig}(L)\) below use this later definition.
\end{counterexample}

The Cont 3 transcription places an
unprimed \(X\) under the supremum over \(X'\ll X\); both the literal
principal-dimension reading and the intended uniform-refinement
reading are stronger than the mature unrestricted theory. The
counterexample below violates even the former.

\begin{definition}[Known: source-resolved optional axioms]
\label{SA21-B-optional-list}
For compatible parents \(F,G\), write \(K=F\cap G\), and take
\(P\ll F,Q\ll G\).
\begin{align*}
 \mathrm C &: P|K\bowtie Q|K\Longrightarrow P\bowtie Q,\\
 \mathrm D &: P|K\parallel Q|K\Longrightarrow P\parallel Q,\\
 \mathrm D_0 &: F\parallel G\Longrightarrow P\parallel Q,\\
 \mathrm L1 &: U_X\ne\varnothing\quad(X\in M),\\
 \mathrm L2 &: X\parallel Y\Longleftrightarrow
                (\forall p\in\ell(X),q\in\ell(Y))\ p\parallel q,\\
 \mathrm L3 &: X\bowtie Y, X\ne Y\Longrightarrow
                (\forall p\in U_X,q\in U_Y)\ p\parallel q.
\end{align*}
The elementary and figure formulations of C and D are equivalent by
the stratum compatibility criterion and carrier restriction.
Chronologically, C is At CD 1 on 156-7/9 and then At C on scan 15;
D is At CD \(1'_0\) on scan 10 and then At D on scan 17;
the hereditary special case is later At D\(_0\).
For arbitrary elementary \(X,Y\), still put \(K=X\cap Y\).
The \emph{absolute} criterion \(\mathrm{CD2}\), also called
\(\mathrm C(\mathrm{spac})\), says that
\[
 (\forall P\ll X,Q\ll Y)
 [P|K=Q|K=0\Longrightarrow P\parallel Q]
 \quad\Longrightarrow\quad X\bowtie Y.              \tag{B.A}
\]
Its polyhedral sufficient version adds \(K=0\) or \(K\in M\) to the
antecedent. The first At CD \(2'\) on scan 11 says only \(K\in M\);
At pol on scan 17 explicitly adds the empty case. Under L2 the empty
case is automatic, so this change does not affect the equivalence below.
The corresponding lieu criterion is
\[
 (\forall p\in\ell(X)\setminus\ell(K),
          q\in\ell(Y)\setminus\ell(K))\ p\parallel q
 \quad\Longrightarrow\quad X\bowtie Y,              \tag{B.CL}
\]
called CL(spac); CL(pol) is the sufficient version with the same extra
condition on \(K\). In contrast, At \(9''\) at 156-6/64 explicitly
states a \emph{biconditional}, and therefore also requires every
compatible pair to have empty or elementary intersection.

At 8, At 9 and At 10 are L1, L2 and L3. At 12 is C.
At 11 is unique gluing over incidence covers; At \(11'\) is its
two-parent, matching-overlap compatibility condition. At 11 bis
concerns filtered unions whose restrictions possess suprema.
At Fil in 156-8 uses the stronger local hypothesis that those
restrictions are eventually stationary.

At 13, or At supp, says that if \(P\ll F\) and some figure is disjoint
from \(P\) but not from \(F\), such a witness can be chosen to refine
\(F\). With L1--L3 and the relative calculus, it makes subdivisions
exactly the incidence-maximal refinements. At 14 is extension of a
prescribed subdivision of a subfigure. At 14 bis adds a cartesian
deployment square, unchanged complementary deployment, and a mixed
incidence rule. At 14 ter asks for a unique elementary extension of
a subdivision of an elementary boundary.
\end{definition}

These are the displayed formulas, not inferred similarities of names:
\path{H156-01_transcription.tex:279--318};
\path{J156-6-08_transcription.tex:141--144,183--198,230--254};
\path{J156-7-01_transcription.tex:8--19,43--57,148--161,182--205};
\path{J156-7-02_transcription.tex:102--143,182--210}; and
\path{J156-8-03_transcription.tex:91--98}.
The original CL(pol) heading on 156-7/23 does not repeat its formula;
our sufficient reading is its stated analogy with At pol on scan 17.
If it is instead meant to include the necessity of empty/principal
intersection, it becomes the stronger At \(9''\)-type condition. We
give the status of both readings wherever it affects a conclusion.

\subsection{Core calculus and exact implications among the optional criteria}

\begin{lemma}[Proved: carrier calculus]
\label{SA21-B-carrier-calculus}
In a prime atelier with At 4--At 7, carriers compose. If \(P\ll F\)
and \(G\leq F\), the displayed formula for \(P|G\) is a figure and
is \(\inf_{\ll}(P,G)\). Restrictions commute with incidence unions
and with incidence intersections inside a fixed parent. If \(F,G\)
are compatible, every common refinement refines \(F\cap G\).
These conclusions require no lieux and no finite-dimension hypothesis.
\end{lemma}

\begin{proof}
If \(Z\ll P\ll F\), let \(A\) be its carrier in \(P\), and let \(B\)
be the carrier of \(A\) in \(F\). Then \(Z\rho A\rho B\), so At 7
gives \(Z\rho B\). Thus \(B\) is the carrier of \(Z\) in \(F\).
The inverse image of a downward-closed deployment under a carrier is
downward closed and is an admitted subfigure. It has exactly the
strata refining \(G\), and it refines \(G\) by At 4.
If \(H\ll P,G\), a stratum \(Z\) of \(H\) has carrier \(A\) in \(P\);
its carrier in \(F\) lies in \(G\). Composition makes that same
stratum the carrier of \(A\) in \(F\), so \(A\ll G\) and
\(Z\ll P|G\). This proves the infimum.
Union and intersection formulas are ordinary inverse-image formulas.
Finally, take the carrier of a common-refining stratum in \(F\vee G\).
Its minimality puts it below a stratum of \(F\) and a stratum of \(G\),
hence in both deployments. At 4 completes the claim.
\end{proof}

\begin{proposition}[Proved: the optional-axiom implication diagram]
\label{SA21-B-optional-implications}
Over the prime core the following implications hold, with exactly the
indicated hypotheses:
\begin{gather*}
 \mathrm C\Longrightarrow\mathrm D\Longrightarrow\mathrm D_0,
 \qquad \mathrm D\Longrightarrow\mathrm L3,
 \qquad \mathrm C\Longrightarrow\mathrm{At}\,11',\\
 \mathrm L2\Longrightarrow\mathrm D_0,
 \qquad \mathrm L2+\mathrm L3\Longrightarrow\mathrm D,
 \qquad \mathrm L3+\mathrm{CL(spac)}\Longrightarrow\mathrm C,\\
 \mathrm D_0+\mathrm L1+\mathrm{CL(spac)}\Longrightarrow\mathrm L2,
 \qquad \mathrm D_0+\mathrm L1+\mathrm{CL(pol)}\Longrightarrow\mathrm L2,\\
 \mathrm L2+\mathrm{CL(spac)}
 \Longleftrightarrow\mathrm L2+\mathrm{CD2},
 \qquad
 \mathrm L2+\mathrm{CL(pol)}
 \Longleftrightarrow\mathrm L2+\mathrm{CD2}' .
\end{gather*}
At \(11'\) is equivalent to finite-cover gluing. Infinite-cover gluing
requires a union-closure hypothesis; At Fil is sufficient. Full
canonical figure families have that closure automatically.
\end{proposition}

\begin{proof}
For C implies D, compatibility follows from C. A common stratum of
\(P,Q\) refines both parents and hence their intersection, so it would
be a common stratum of the disjoint restrictions. D implies D\(_0\)
by taking empty intersection. To obtain L3 from D, take the singleton
figures \(p\rho X,q\rho Y\) for distinct compatible \(X,Y\).
At least one restricts to zero on \(X\cap Y\): both could be nonzero
only if both \(X,Y\) belonged to that intersection. Their restrictions
are therefore disjoint. D yields \(p\parallel q\). Matching
restrictions are compatible, so C implies At \(11'\).

L2 immediately gives hereditary disjointness. For L2+L3 implies D,
choose \(p\ll P,q\ll Q\). If both refine \(K\), use L2 for the
disjoint restrictions. Otherwise their carriers in \(F\vee G\) are
distinct, and L3 applies. L2 then gives \(P\parallel Q\).
For L3+CL(spac) implies C, it suffices by the stratum criterion
to take elementary \(P,Q\): use the carriers of their strata as
elementary parents, and then test every cross-pair. Take places outside
\(\ell(P\cap Q)\). If both refine \(K\), their carriers in the
compatible union \((P|K)\vee(Q|K)\) are distinct; if one does not,
their carriers in \(F\vee G\) are distinct. In either case L3 gives
orthogonality, and CL(spac) gives compatibility.

For the next two arrows, hereditary disjointness proves the forward
half of L2. If all places of \(X\) and \(Y\) are pairwise disjoint,
L1 forces \(X\cap Y=0\), since any nonzero intersection would contain
a place orthogonal to itself. Either CL criterion then gives
compatibility, and hence disjointness.

For the absolute equivalences, assume L2. The hypothesis in (B.A)
implies the hypothesis in (B.CL), by testing singleton places outside
\(\ell(K)\). Conversely, if (B.CL)'s hypothesis holds and
\(P|K=Q|K=0\), every place of \(P\) or \(Q\) is outside \(\ell(K)\),
by the infimum property. L2 gives \(P\parallel Q\), proving (B.A)'s
hypothesis. The empty/principal qualification is unchanged throughout.

For finite gluing, matching restrictions make the local pieces
pairwise compatible by At \(11'\); their finite union is a figure,
refines the covered parent, and has exactly the prescribed
restrictions. Any other gluing is their union because every stratum
has a carrier in some cover member. Conversely apply gluing to the
cover by two compatible parents. For an arbitrary cover, take the
increasing family of gluings over finite subcovers. On each stratum
of the parent it is eventually equal to the prescribed local
refinement. At Fil supplies its union. The same argument works
without At Fil if every clique downset is admitted.
\end{proof}

The finite-cover assertion does not imply L3 by itself: nested parents
do not furnish two independent cover pieces. The manuscript explicitly
singles out that nested case
\citep[156-6, scan 65, marginal note]{Grothendieck1986},
\path{J156-6-08_transcription.tex:248--250}.
We provide an explicit separating model below.

\subsection{The magasin formulations and the corrected scope of L6}

\begin{theorem}[Proved: the strong relational equivalences]
\label{SA21-B-mag-equivalence}
The following two structures are equivalent, by mutually inverse
definitions on the same set \(M\).
\begin{enumerate}
\item Mag 1--4: \(\leq,\ll\) are partial orders,
\(\leq\subseteq\ll\); every nonempty
\(\{Z\leq Y:X\ll Z\}\) has a least element; \(\delta\) is symmetric,
irreflexive and hereditary under the resulting \(\rho\); and
\[
 Y,Z\leq X,\ Y\ne Z\quad\Longrightarrow\quad Y\delta Z.
                                                        \tag{B.M4}
\]
\item The corrected strong Mag \(1'\)--\(4'\): \(\leq,\rho\) are
partial orders; \(\delta\) is symmetric and irreflexive; there is at
most one \(Z\) with \(X\rho Z\leq Y\); the exchange condition
\[
 X'\leq X\rho Y\quad\Longrightarrow\quad
             \exists Y'\leq Y\ (X'\rho Y')             \tag{B.E}
\]
holds; \(\delta\) is hereditary under \(\rho\); and (B.M4) holds.
\end{enumerate}
The inverse definition is
\(X\ll Y\Longleftrightarrow\exists Z\ (X\rho Z\leq Y)\).

For either structure define
\[
 X\bowtie Y\Longleftrightarrow
 \bigl(\forall A\in\widetilde X\setminus\widetilde Y,
             B\in\widetilde Y\setminus\widetilde X\bigr)\ A\delta B.
                                                        \tag{B.MC}
\]
Every principal downset is a clique. The full family of clique downsets,
or any atlas containing all principal downsets and closed under
subfigures and finite compatible unions, is a prime atelier with
At 1--7, C, D, D\(_0\), L3 and finite gluing. For the full family,
arbitrary gluing, At 11 bis and At Fil also hold.
\end{theorem}

\begin{proof}
Start with the primed data. Reflexivity and \(\leq\subseteq\ll\) are
immediate. If \(X\rho A\leq Y\rho B\leq Z\), exchange supplies
\(A\rho C\leq B\); transitivity of \(\rho\) gives \(X\rho C\leq Z\).
Thus \(\ll\) is transitive. If also \(Y\ll X\), write
\(X\rho A\leq Y\rho B\leq X\). Exchange again gives
\(A\rho C\leq B\leq X\). The uniqueness clause for \((X,X)\)
forces \(C=X\), since \(X\rho X\leq X\). Antisymmetry of \(\rho\)
then gives \(A=X\), hence \(X\leq Y\). Uniqueness for \((Y,Y)\)
similarly gives \(B=Y\), hence \(Y\leq X\). Therefore \(X=Y\).
If \(X\rho A\leq Y\) and \(X\ll Z\leq Y\), factorization through
\(Z\) and uniqueness force the same \(A\), so \(A\leq Z\).
Consequently \(A\) is the least carrier, and the recovered internal
refinement is exactly the original \(\rho\).

Conversely, carrier existence gives unique factorization, and applying
it to \(X'\leq X\rho Y\) gives exchange. To prove internal
transitivity, suppose \(X\rho Y\rho Z\), and let
\(X\rho W\leq Z\) be the least-carrier factorization. If \(W<Z\),
(B.M4) gives \(W\delta Z\). Heredity first along \(X\rho W\), then
along \(Y\rho Z\), then along \(X\rho Y\), gives \(X\delta X\),
contrary to irreflexivity. Thus \(W=Z\), proving transitivity.

For the figure construction, (B.M4) and (B.MC) imply that every pair of
distinct strata in a compatible union is \(\delta\)-related: pairs
within a principal deployment use (B.M4), while the remaining cross
pairs use (B.MC). Compatibility is hereditary under incidence.
If a stratum \(Z\) refines two strata of one figure, factor it
internally through carriers \(A,B\). Distinct \(A,B\) would be
\(\delta\)-related and heredity would give \(Z\delta Z\).
Thus the carriers coincide, proving At 5. If compatible \(X,Y\)
satisfy \(X\ll Y\), factor \(X\rho A\leq Y\). If \(A\ne X\),
compatibility gives \(X\delta A\), and heredity gives \(X\delta X\).
Hence \(X=A\leq Y\), proving At 6.

Refinement of figures is defined by At 4. Its antisymmetry for
\emph{arbitrary}, possibly infinite figures follows from carriers:
if \(F\ll G\ll F\), the successive carriers of \(X\in F\) are
\(Y\in G,Z\in F\), with \(X\rho Y\rho Z\). But the carrier of
\(X\) in \(F\) is \(X\), so \(Z=X\); antisymmetry on \(M\) gives
\(Y=X\). Thus \(F\subseteq G\), and symmetrically equality.
Generation and primality are union statements for downsets.

To prove C, let \(F,G\) be compatible parents and let
\(P\ll F,Q\ll G\) have compatible restrictions to \(K=F\cap G\).
For distinct \(A\in P,B\in Q\), take their carriers in the parent
union. Distinct carriers are internally disjoint, hence \(A\delta B\).
If the carriers agree, they lie in \(K\), so both strata occur in the
compatible restrictions and are again internally disjoint. Thus the
union is a clique downset, and the atlas's finite-union condition
admits it. D, D\(_0\) and gluing follow from
Proposition~\ref{SA21-B-optional-implications}. L3 follows directly
from heredity of \(\delta\) along places. In the full family every
increasing compatible union is admitted; the infinite assertions
follow as well.
\end{proof}

\begin{theorem}[Proved: the precise local equivalence]
\label{SA21-B-local-equivalence}
For a strong magasin, Mog L2 is
\[
 X\delta Y\Longleftrightarrow
       (\forall p\in U_X,q\in U_Y)\ p\delta q.        \tag{B.Loc}
\]
It implies Mog L1, \(U_X\ne\varnothing\), by irreflexivity. Under the
prime incidence convention, strong magasins with (B.Loc) and their
chosen atlases are equivalent to prime ateliers satisfying
At 1--7, L1, L2, L3 and CL(spac). The inverse definition of
\(\delta\) is the right-hand side of (B.Loc).

In the earlier formulation using \(\leq,\ll\), Mog 1, Mog 2,
Mog L1 and a symmetric irreflexive graph on the lieux, together with
the all-pairs div 4 clause, produce a strong magasin by this definition
\emph{if and only if the resulting internal refinement is transitive}.
It suffices to impose transitivity when the first object is a lieu;
under these hypotheses that restricted condition forces full transitivity.
In the later \(\leq,\rho\) formulation, the exact repair is to require
nonempty \(U_X\) and replace comparable-only div 3 by the all-pairs
condition (B.M4) tested on open lieux.
\end{theorem}

\begin{proof}
The forward half of (B.Loc) is Mag 3; its converse is Mog L2.
If \(U_X=\varnothing\), that converse would give \(X\delta X\).
The strong magasin already supplies L3. Its disjointness formula is
\[
 F\parallel G\Longleftrightarrow
       (\forall X\in\widetilde F,Y\in\widetilde G)\ X\delta Y.
                                                        \tag{B.MD}
\]
Indeed the right side ensures compatibility by (B.MC) and rules out
a shared stratum; the converse uses (B.M4) in the compatible union.
Decompose each lieu by its carrier and use (B.Loc) in (B.MD) to obtain
L2. For CL(spac), a stratum \(A\) of \(X\) absent from \(Y\) has
all its open lieux outside \(\ell(X\cap Y)\), by carrier uniqueness.
The exterior-lieu criterion therefore gives \(A\delta B\) for every
cross-pair in (B.MC), proving compatibility.

Conversely define \(\delta\) by open-lieu orthogonality. L1 gives
irreflexivity; symmetry is immediate; transitivity of internal
refinement gives \(U_A\subseteq U_X\) when \(A\rho X\), hence
heredity. L3 gives (B.M4). To recover the original compatibility,
its necessity follows from L3. For sufficiency in (B.MC), take
places outside the common figure and use their distinct exterior
carriers; their \(\delta\)-relation gives the hypothesis of CL(spac).
This recovers the original structure and proves the equivalence.

In the earlier data, internal transitivity makes the same heredity
argument work, while the all-pairs div 4 clause gives (B.M4).
The heredity argument only uses \(p\rho A\rho X\Rightarrow p\rho X\)
for lieux \(p\). Thus this restricted transitivity already constructs
a strong magasin, whose full internal transitivity then follows from
Theorem~\ref{SA21-B-mag-equivalence}.
Necessity follows from Theorem~\ref{SA21-B-mag-equivalence}, which
derives internal transitivity in every strong magasin. In the later
data that transitivity is primitive; nonempty open-lieu sets are
exactly what makes the reconstructed relation irreflexive, and
the all-pairs clause is exactly (B.M4).
\end{proof}

The source locators are
\path{J156-8-01_transcription.tex:243--263} for strong Mag 1--4;
\path{J156-8-02_transcription.tex:285--315,330--351} for the open
lieux and Mog L1--L2; and
\path{J156-8-04_transcription.tex:199--247,280--300} for the primed
and later divisible presentations. The relation in the final
Mag \(4'\) line carries a transcription doubt before \(X\mid0\mid Y\).
The original scan 31 was inspected at enlarged resolution: a separate
glyph is visible there, but its reading as a negation remains uncertain.
The intended positive comparable-pair reading is already too weak;
the negative reading is incompatible with the earlier positive
two-strata rule whenever there is a strict incidence. Neither reading
can be silently identified with (B.M4).

\begin{counterexample}[Refuted: the missing transitivity in 156-8/16--17]
\label{SA21-B-missing-transitivity}
Take \(M=\{p,q,z,u,t\}\), incidence generated by \(z<t\), and
refinement generated by
\[
 p\ll z,\quad p\ll u,\quad q\ll u,\quad u\ll t,
 \qquad z\ll t.
\]
Put \(p\parallel q\) on the two lieux. The orders contain identity
pairs. Least carriers exist in every principal deployment. The open
lieu sets are
\[
 U_p=U_z=\{p\},\quad U_q=U_t=\{q\},\quad U_u=\{p,q\}.
\]
Thus Mog 1, Mog 2, divisibility, symmetric irreflexive div 3, and
the all-pairs div 4 clause all hold: the only nontrivial incidence
pair is \(z<t\), with disjoint open lieu sets.
But \(p\rho u\rho t\), whereas \(p\not\rho t\), since \(p\ll z<t\).
The defined relation has \(p\delta t\) but not \(p\delta u\), violating
Mag 3 along \(u\rho t\). This independently verifies the input
counterexample SA21-R8-09 and establishes the exact missing condition
in Theorem~\ref{SA21-B-local-equivalence}.
\end{counterexample}

\begin{counterexample}[Refuted: the weak primed and divisible lists]
\label{SA21-B-weak-primed}
For the weak primed list take \(M=\{y,z,t\}\),
\(y<t,z<t\), \(\rho=\mathrm{id}\), and
\(\delta=\{(y,t),(t,y),(z,t),(t,z)\}\).
The two orders, unique factorization, exchange and heredity all hold.
Every strict comparable pair is internally disjoint. But \(y,z<t\)
are not internally disjoint, so the supposed principal deployment
\(\{y,z,t\}\) is not a figure under (B.MC).

If the doubtful mark in final Mag \(4'\) is read as a negation,
take instead the two-element chain \(a<t\), with
\(\rho=\mathrm{id}\) and \(\delta=\varnothing\).
The two partial orders, unique factorization, exchange, symmetry,
irreflexivity and heredity all hold, as does the negative
comparable-pair clause. Strong Mag 4 fails at \(a<t\).
Thus this alternative reading does not restore the equivalence;
the certificate checks its own full primed hypothesis list.

For the later divisible list, including explicit divisibility, take
\(M=\{a,b,c,t\}\), incidence \(a<t,b<t\), and internal refinement
\(c\rho t\), besides identities. The lieux are \(a,b,c\), with graph
edges \(a\parallel c,b\parallel c\), and no edge \(a\parallel b\).
Unique factorization and exchange hold; every open-lieu set is a
singleton. The comparable div 3 clause holds for \(a<t,b<t\).
The derived \(\delta\) has edges
\(ac,bc,at,bt\), but not \(ab\). Again the principal deployment
of \(t\) is not a figure. These are genuine finite countermodels to
the written weakened clauses, not failures of the corrected clauses.
\end{counterexample}

\begin{counterexample}[Refuted: L6, including strengthened nonimplications]
\label{SA21-B-L6-countermodels}
All the following strong magasins use the full canonical figure
family. Relations not generated by the displayed pairs are absent.
\begin{enumerate}
\item Let incidence be equality on \(\{p,q,x,y\}\), refinement be
generated by \(p\ll x,q\ll y\), and let the sole undirected
\(\delta\)-edge be \(pq\). Mag 1--4 and divisibility hold.
Nevertheless \(U_x=\{p\},U_y=\{q\}\) are orthogonal while
\(x\not\delta y\); L2 and CL(spac) fail. Every class \(\Sigma\)
containing the empty poset and the occurring principal posets also
fails At \(9_\Sigma\) on \(x,y\). Thus L6 remains false even if
its excluded divisibility axiom is supplied.
\item Let incidence and refinement both consist of the two chains
\(a<x,b<y\). Let \(\delta\) have edges \(ax,by,ab\).
Internal refinement is equality, so Mag 1--4 hold. The only lieux
below \(x,y\) are respectively \(a,b\), which are disjoint, but
\(x,y\) are not disjoint. This is a second direct failure of L2.
\item On \(\{a,x,y\}\), let incidence and refinement be
\(a<x,a<y\), and let \(\delta\) have just the edges \(ax,ay\).
Mag 1--4, L2, L3 and At supp hold. The two shadows of \(x,y\)
are both \(\{a\}\), as is the shadow of their common incidence
figure. The exterior criterion is empty and hence true, but
\(x,y\) are incompatible. Thus CL(spac), CL(pol), CD2 and
\(\mathrm{CD2}'\) all fail, even with At supp. For CD2, every nonempty refinement
has a nonempty restriction to \(a\); its antecedent is vacuous.
\item On \(\{a,b,c,d\}\), take incidence and refinement
\(a<b<c\), with \(d\) isolated, and let \(\delta\) have edges
\(ab,ac,bc,ad,bd\), omitting \(cd\). Mag 1--4 hold.
The seven figures are
\[
 0,\ a,\ d,\ a+d,\ a+b,\ a+b+d,\ a+b+c.
\]
The principal figure \(a\) subdivides \(b\), because their disjoint
figures are exactly \(0,d\). The only subdivision of \(c\) is \(c\)
itself, whose restriction to \(b\) is \(b\). Therefore this
subdivision of a subfigure does not extend: At 14 fails.
\end{enumerate}
Every pair, triple, figure and refinement in these models is checked
by \texttt{core\_models.py}; the proofs of the general magasin
construction verify the incidence and figure axioms as well.
\end{counterexample}

\begin{counterexample}[Refuted: finite gluing or D\(_0\) need not give D]
\label{SA21-B-gluing-not-L3}
Take \(M=\{x,y,q\}\), incidence \(x<y\), refinement
\(x\ll y,q\ll y\), and compatibility consisting of equality and
the pair \(x,y\). The figures are
\(0,\{x\},\{q\},\{x,y\}\). The lieux are \(x,q\), and
\(U_y=\{q\}\). The prime core, L1, L2, CL(spac), At supp and
D\(_0\) hold. No two nonempty figures are disjoint, so L2 and
D\(_0\) have exactly their stated, mostly vacuous content.
Every incidence cover of \(\{x,y\}\) contains \(\{x,y\}\) itself;
the other nonzero figures are elementary singletons. Consequently
all compatible local refinement data glue uniquely. This checks
At 11, including arbitrary covers after deleting repetitions.

For parents \(x\leq y\), the refinements \(x\ll x\) and \(q\ll y\)
restrict respectively to \(x\) and \(0\) on their intersection.
Those restrictions are disjoint, while \(x,q\) are not. Thus D
and L3 fail, and so does C. The exhaustive code verifies every
finite cover and all its refinement assignments. This also shows
why CL(spac) alone must not be confused with the relative C axiom.
\end{counterexample}

\begin{counterexample}[Refuted: At Fil for a restricted atlas]
\label{SA21-B-fil}
Let \(M=\{a_n,x_n:n\in\mathbb N\}\), with discrete incidence,
refinement generated by \(a_n\ll x_n\), and internal disjointness
between exactly those pairs having different indices. This is a
strong local magasin: the open-lieu set of either member of the
\(n\)-th pair is \(\{a_n\}\). A full figure chooses at most one
member of each pair. Restrict the atlas to those figures containing
only finitely many \(a_n\), allowing arbitrarily many \(x_n\).
It contains all principal figures and is closed under subfigures,
bounded unions within the atlas and finite compatible unions.
It therefore satisfies all the prime core, local and finite-gluing
axioms proved above.

Let \(F=\{x_n:n\in\mathbb N\}\) and
\(G_k=\{a_0,\ldots,a_k\}\). These form an increasing family of
refinements of \(F\); their restriction to each stratum \(x_n\)
is eventually constant. Their union \(\{a_n:n\in\mathbb N\}\)
is omitted, and no admitted figure can bound it. Thus At Fil,
At 11 bis and arbitrary gluing fail. For the last assertion use the
disjoint cover of \(F\) by the singleton figures \(x_n\), with local
refinements \(a_n\). This identifies the exact extra atlas closure
needed by the otherwise valid gluing theorem.
\end{counterexample}

\begin{proposition}[Proved: extension after the two missing hypotheses]
\label{SA21-B-extension-Zorn}
For the full figure family of a strong magasin satisfying L1 and
At supp, subdivision is equivalent to incidence-maximal refinement,
and At 14 follows. This assertion uses the usual axiom of choice
through Zorn's lemma. It does not give uniqueness in At 14 bis or ter.
\end{proposition}

\begin{proof}
Let \(P\preceq F\) and \(P<Q\ll F\). Choose a stratum
\(X\in Q\setminus P\) and a place \(p\rho X\), using L1.
All strata of \(P\) are distinct from \(X\) in the compatible
figure \(Q\). Mag 3--4 and (B.MD) give \(p\parallel P\).
But \(p\ll F\), so \(p\) cannot be disjoint from \(F\): D\(_0\)
would give \(p\parallel p\). This contradicts equal cosupports.
Thus a subdivision is maximal. Conversely, if a maximal refinement
\(P\ll F\) has a different cosupport, At supp supplies a nonzero
\(H\ll F\) disjoint from \(P\), and \(P\cup H\) is a strictly
larger refinement. This proves the equivalence.

Now fix \(G\leq F\) and \(G'\preceq G\). Consider refinements
\(P\ll F\) with \(G'\leq P\) and \(P|G=G'\). The figure \(G'\)
itself belongs to this collection. Every chain has its union as
an upper bound: it is a clique downset, refines \(F\), and its
restriction remains \(G'\). Let \(P\) be maximal in this collection.
Any larger refinement \(Q\ll F\) would have
\(G'\leq Q|G\ll G\). The already proved maximality of the subdivision
\(G'\) forces \(Q|G=G'\), contradicting the choice of \(P\).
Thus \(P\) is a maximal refinement of \(F\), hence a subdivision.
\end{proof}

\begin{remark}[Proved and refuted: the exact replacement for L6]
\label{SA21-B-L6-status}
For strong Mag 1--4, the automatic relative package is the prime core,
C, D, D\(_0\), L3 and finite gluing; arbitrary gluing and At Fil are
automatic for the full canonical figure family. The local package
L1, L2 and CL(spac) requires Mog L2, or equivalently the hypotheses
of Theorem~\ref{SA21-B-local-equivalence}. Extension requires additional
hypotheses, as Proposition~\ref{SA21-B-extension-Zorn} specifies.
None of finite dimension, polyhedrality, simpliciality, incidence C4,
or unique cone extension follows from Mag 1--4. For finite dimension
take any infinite-height incidence poset, \(\rho=\mathrm{id}\) and
\(\delta\) unequal-pair disjointness. Even local full set-theoretic
ateliers can have infinite-height principal deployments: on
\(\mathbb N\cup\{\infty\}\), use the elementary admissible family
of finite initial segments together with the whole set. The
polyhedral and simplicial failures are proved in Part C, and C4 is
settled in Part B.2 below. Therefore the proposed exception list in
L6 is false, and the replacements above specify both the valid
consequences and the necessary qualifications.
\end{remark}

\subsection{\texorpdfstring{At \(9_\Sigma\)}{At 9(Sigma)} plus At 10: a counterexample and the exact obstruction}

\begin{theorem}[Proved: the exact remaining \(\Sigma\)-condition]
\label{SA21-B-Sigma-gate}
Assume the prime core and At 10. Let \(\Sigma\) be a replete class of
posets containing all occurring principal deployments. Suppose
At \(9_\Sigma\) states
\[
 X\bowtie Y\Longleftrightarrow
 \left\{
 \begin{array}{l}
 (\forall p\in\ell(X)\setminus\ell(K),
            q\in\ell(Y)\setminus\ell(K))\ p\parallel q,\\
 \widetilde K\in\Sigma,
 \end{array}\right.
 \qquad K=X\cap Y.                                  \tag{B.Sigma}
\]
For compatible parents \(X,Y\), elementary refinements
\(X'\ll X,Y'\ll Y\), and compatible restrictions to \(K=X\cap Y\),
the exterior-lieu condition in (B.Sigma) for \(X',Y'\) is automatic.
Consequently
\[
 X'\bowtie Y'
 \quad\Longleftrightarrow\quad
 \widetilde{X'\cap Y'}\in\Sigma.                     \tag{B.Gate}
\]
Thus At 12 follows exactly when these common-refinement deployments
remain in \(\Sigma\). In particular it follows for unrestricted
\(\Sigma\). For \(\Sigma\) equal to the empty poset and posets with
a greatest element, it follows if elementary refinements restrict
to empty or elementary figures on every elementary subfigure.
\end{theorem}

\begin{proof}
Put \(A=X'|K,B=Y'|K\), and \(K'=X'\cap Y'\). The common-refinement
lemma gives \(K'=A\cap B\). Choose
\(p\in\ell(X')\setminus\ell(K')\) and
\(q\in\ell(Y')\setminus\ell(K')\).
If both refine \(K\), their carriers in \(A\vee B\) belong
respectively to \(A,B\). They cannot coincide: a common carrier
would put \(p\) in \(\ell(K')\). At 10 gives \(p\parallel q\).
If one, say \(p\), does not refine \(K\), its carrier in
\(X\vee Y\) lies in \(\widetilde X\setminus\widetilde Y\), whereas
the carrier of \(q\) lies in \(\widetilde Y\); they are distinct,
and At 10 again applies. This proves the automatic first clause
and hence (B.Gate).

For the final assertion, compatibility of the parents makes
\(K\) empty or elementary. The added restriction hypothesis makes
\(A,B\) empty or elementary. If either is empty, so is \(K'\).
Otherwise compatibility of the elementary \(A,B\), through
(B.Sigma), puts their common deployment in \(\Sigma\).
\end{proof}

\begin{counterexample}[Refuted: At \(9_\Sigma\)+At 10 does not imply At 12]
\label{SA21-B-Sigma-counter}
Take the underlying set \(\{a,b,c,p\}\). We use seven elementary
admissible families, with point objects denoted by their points:
\[
 \begin{array}{c|l}
 a,b,c,p & \{\{a\}\},\{\{b\}\},\{\{c\}\},\{\{p\}\}\\
 k & \{\{a\},\{b\},\{a,b,c\}\}\\
 x & \{\{a\},\{b\},\{a,b,c\},\{a,b,c,p\}\}\\
 x' & \{\{a\},\{b\},\{a,b,p\}\}.
 \end{array}
\]
Incidence is closed-subfamily inclusion; refinement is ordinary
set-figure refinement. Declare a pair compatible precisely when
it is set-theoretically compatible and its common deployment is
empty or elementary. The figures are all downward-closed cliques
for this compatibility. The complete nontrivial data are
\[
\begin{array}{c|c|c|c}
 Z&\widetilde Z&\{W:W\ll Z\}&\{W:W\rho Z\}\\ \hline
 a&\{a\}&\{a\}&\{a\}\\
 b&\{b\}&\{b\}&\{b\}\\
 c&\{c\}&\{c\}&\{c\}\\
 p&\{p\}&\{p\}&\{p\}\\
 k&\{a,b,k\}&\{a,b,c,k\}&\{c,k\}\\
 x&\{a,b,k,x\}&\{a,b,c,p,k,x,x'\}&\{p,x,x'\}\\
 x'&\{a,b,x'\}&\{a,b,p,x'\}&\{p,x'\}.
\end{array}
\]
The incompatible distinct pairs are exactly
\[
 (c,k),(c,x),(p,x),(p,x'),(k,x'),(x,x').
\]
There are 21 figures: the sixteen subsets of the four point objects,
and
\[
 \{a,b,k\},\quad\{a,b,k,p\},\quad\{a,b,k,x\},\quad
 \{a,b,x'\},\quad\{a,b,c,x'\}.
\]
Take \(\Sigma\) to be the class consisting of the empty poset and
all posets having a greatest element.

Every principal deployment is a figure. The table gives two partial
orders, unique carriers for every stratum refining every figure,
coincidence of incidence and refinement on compatible objects, and
transitive internal refinement. Thus the prime At 1--7 hold.
The lieux are exactly \(a,b,c,p\), and the open lieu sets of
\(k,x,x'\) are respectively \(\{c\},\{p\},\{p\}\); hence At 8.
Distinct points are disjoint, and ordinary support disjointness is
exactly figure disjointness, giving At 9. Compatible distinct
objects have disjoint open lieu sets, giving At 10. The explicit
compatibility definition is exactly (B.Sigma), so At \(9_\Sigma\)
holds. These assertions, including all figure carriers, are
exhaustively checked in \texttt{finite\_axioms.py}.

Now use compatible parents \(k\leq x\), with refinements
\(k\ll k\) and \(x'\ll x\). Their restrictions to \(K=k\) are
\(k\) and \(\{a,b\}\), which are compatible. But \(k,x'\) are
incompatible: their common deployment \(\{a,b\}\) has two incomparable
maximal elements and is not in \(\Sigma\). This refutes the proposed
implication. The same program verifies D and the matching-overlap
criterion for \emph{every} compatible figure-parent pair and every
refining figure pair. Therefore At 11 also holds, by finite gluing.
Adding At 11 and At D does not repair the implication.
\end{counterexample}

The source question is
\path{J156-6-10_transcription.tex:122--145}, scans 78--80,
and the diagram summarized in input check SA21-R6-37. We independently
confirm its empty-poset qualification: if \(\varnothing\notin\Sigma\),
even two disjoint point figures can never be compatible under
(B.Sigma). With L1 and \(\varnothing\in\Sigma\), (B.Sigma) implies
L2: orthogonal whole shadows force an empty common figure, and the
biconditional then gives compatibility; its necessity gives the
reverse direction. The positive unrestricted result is the precise
At \(9'\)+At 10 implication, while (B.Gate) identifies the exact
additional closure needed in every other \(\Sigma\)-case.

\subsection{At 14 bis: deployment uniqueness and object uniqueness}

\begin{counterexample}[Refuted: unique extension in At 14 bis and ter]
\label{SA21-B-extension-unique}
Let \(M=\{a,x\}\), let incidence be equality, let refinement have
\(a\ll x\), and let internal disjointness be empty. The full figures
are \(0,a,x\); two nonempty figures are compatible only when equal,
and no nonempty figures are disjoint. Mag 1--4 hold: every principal
deployment is a singleton, carriers are unique, and all orthogonality
clauses are vacuous. The sole lieu is \(a\), and
\(U_a=U_x=\{a\}\). Thus Mog L2, the local axioms, the prime core,
C, D, all gluing and At supp hold. Every nonempty refinement has
the same cosupport, so \(a\preceq x\).

Take \(F=x\), \(G=G'=0\). Both choices \(F'=x\) and \(F'=a\)
are subdivisions. Their deployment maps are isomorphisms of
one-point posets, their restrictions over \(G\) are empty, and the
deployment squares are cartesian. The complementary deployment
map is an isomorphism. The mixed incidence condition is vacuous.
Hence both satisfy At 14 bis, but they are different figures.
Their boundaries are both empty, so At 14 ter's uniqueness also
fails. This model satisfies At 14 bis \emph{existence for every
possible diagram}: the finite certificate enumerates all six
extension diagrams and all twelve gluing assignments.
\end{counterexample}

\begin{proposition}[Proved: the exact uniqueness supplied by the diagram]
\label{SA21-B-deployment-unique}
Fix posets \(\widetilde G\subseteq\widetilde F\), with
\(\widetilde G\) a downset, and a monotone map
\(r:\widetilde G'\to\widetilde G\). A deployment square with
inverse image \(\widetilde G'\), an isomorphism over
\(\widetilde F\setminus\widetilde G\), and the intended mixed
incidence rule is unique up to unique isomorphism over these data.
Its poset is
\[
 E=\widetilde G'\amalg(\widetilde F\setminus\widetilde G),
\]
with the original orders on each piece, no exterior-to-interior
relations, and
\[
 \xi\leq\eta\quad\Longleftrightarrow\quad
 r(\xi)\leq\eta
 \qquad(\xi\in\widetilde G',\ 
        \eta\in\widetilde F\setminus\widetilde G).
\]
This is a uniqueness theorem for an abstract ordered deployment.
An atelier can have several objects with this same deployment and
carrier map, as Counterexample~\ref{SA21-B-extension-unique} shows.
Requiring a coarsest extension is an additional property, not a
consequence of the displayed deployment diagram.
\end{proposition}

\begin{proof}
Every element belongs either to the prescribed inverse image or to
the identified complementary poset. The within-piece orders and
mixed rule prescribe every possible ordered pair. There are no
relations from the exterior to the interior because the latter is
a downset. The displayed relation is transitive: the only mixed
transitivity checks use monotonicity of \(r\) or transitivity in
\(\widetilde F\). Antisymmetry follows because no reverse mixed
pair exists. Thus it is a poset, and the prescribed identifications
give the unique isomorphism between any two such diagrams.
\end{proof}

We inspected the actual scan 156-6/88, not only its transcription.
The written mixed clause quantifies \(\eta\in\widetilde F'\), not
\(\eta\in\widetilde F'\setminus\widetilde G'\):
\path{J156-6-12_transcription.tex:4--24}.
Read literally, it also requires the map
\(\widetilde G'\to\widetilde G\) to reflect order. A genuine
subdivision with several incomparable strata carried by the same
old stratum violates that requirement. Part B.4 gives such a closed
interval subdivision and its complete finite carrier check.
The geometric mixed-boundary reading gives
Proposition~\ref{SA21-B-deployment-unique}; the literal stronger
reading can fail even in that geometric model. The two-object
counterexample works under \emph{both} readings because
\(\widetilde G'=\varnothing\).

\subsection{The converse of the 156-6 Th\'eor\`eme--Scholie}

\begin{theorem}[Proved: the corrected representation, in both directions]
\label{SA21-B-Atstar}
The groupoid of prime ateliers with At 4--At 7 is equivalent to the
groupoid of pairs \((E,\mathscr H\subseteq\operatorname{Fig}(E))\)
with the following five properties:
\begin{enumerate}
\item \(\varnothing\in\mathscr H\), and \(\mathscr H\) is closed
under unions bounded by a member of \(\mathscr H\).
\item If \(A\in\Phi\cap\Psi\), the principal subfigures
\(\Phi_A=\{B\in\Phi:B\subseteq A\}\) and \(\Psi_A\) agree,
\emph{and their common value belongs to \(\mathscr H\)}.
Denote it by \(\widehat A\).
\item For occurring strata \(A\subseteq B\),
\(\widehat A\ll\widehat B\) as set-theoretic figures.
\item If \(\widehat A\cup\widehat B\in\mathscr H\) for every
\(A\in\Phi,B\in\Psi\), then \(\Phi\cup\Psi\in\mathscr H\).
\item \(\mathscr U=\bigcup_{\Phi\in\mathscr H}\Phi\) is an
admissible family of support \(E\), and it separates the points of \(E\).
\end{enumerate}
The italicized membership clause is present at scan 32, but omitted
from the shortened final At\(^*2\) list at scan 34. With that clause
retained, the converse is valid. Without it, if all displayed
refinements are interpreted in the ambient \(\operatorname{Fig}(E)\),
the final list is insufficient. Without (B.P) on the atelier side,
the forward equivalence is already false by
Counterexample~\ref{SA21-B-M3}.
\end{theorem}

\begin{proof}
First start from \((E,\mathscr H)\). Every point \(e\in E\) has a
least member \(A_e\in\mathscr U\) containing it, by admissibility.
The residual of each \(A\in\mathscr U\) is nonempty. If two points
have the same least member, they belong to exactly the same members
of \(\mathscr U\); separation makes them equal. Thus every residual
is a singleton \(e_A\), and \(e\mapsto A_e\) is a bijection
\(E\to\mathscr U\).

Transfer inclusion of these sets to a refinement order on \(E\):
\(e_A\ll e_B\) iff \(A\subseteq B\). Transfer occurrence in
canonical principal figures to incidence:
\(e_A\leq e_B\) iff \(A\in\widehat B\). This is a partial order:
membership in a principal subfigure is transitive, and antisymmetry
follows from set inclusion. Property 2 makes every
\(\Phi\in\mathscr H\) a downset for this incidence. Properties 1
and 2 admit every closed subfamily, since it is the bounded union
of its principal subfigures. The principal members are exactly the
completely join-prime irreducibles, and property 4 is compatibility
by strata. This establishes the prime incidence core.

For a point \(e_A\) refining \(\Phi\), the least member of \(\Phi\)
containing \(e_A\) is the least member containing \(A\):
\(A\) is already the least member of the entire \(\mathscr U\)
containing that point. This gives the least-carrier axiom.
If \(\widehat A,\widehat B\) are compatible and \(A\subseteq B\),
their admitted union has \(A\) as a substratum of \(B\);
property 2 forces \(A\in\widehat B\). Thus compatible refinement
is incidence.

Refinement of two members \(\Phi,\Psi\) is equivalent to every
\(A\in\Phi\) being contained in some \(B\in\Psi\). Necessity is
the monotone-carrier property of set-figure refinement. For
sufficiency, property 3 gives
\(\widehat A\ll\widehat B\leq\Psi\). The open stratum of \(A\)
in \(\Phi\) is its open stratum in \(\widehat A\), so these
refinements supply a unique target cell for every source cell.
The carrier maps agree on principal subfigures and are monotone;
therefore \(\Phi\ll\Psi\). This is At 4. If
\(e_A\rho e_B\rho e_C\), the corresponding principal
refinements carry the top open cell first into the top cell of
\(\widehat B\), then into that of \(\widehat C\). Composition of
cell containment gives \(e_A\rho e_C\), proving At 7.

Conversely start from a prime atelier. Set \(E=M\),
\(A_X=\{Z\in M:Z\ll X\}\), and
\(\Phi_F=\{A_X:X\in\widetilde F\}\).
The principal \(\ll\)-ideals \(A_X\) form an admissible separating
cover of \(M\), with residual singleton \(\{X\}\), giving property 5.
Least carriers make every \(\Phi_F\) admissible. For strata inside
one figure, \(A_X\subseteq A_Y\) means \(X\ll Y\) and hence
\(X\leq Y\); therefore the canonical subfigure at \(A_X\) is
exactly \(\Phi_X\), giving property 2. Bounded union preservation
and stratum compatibility give properties 1 and 4.
Carrier composition gives set-figure refinement
\(\Phi_X\ll\Phi_Y\) whenever \(X\ll Y\), hence property 3.
The point-to-least-stratum bijection constructed in the first half
recovers \(X\mapsto A_X\); the figures, incidence and refinement
are recovered exactly. Isomorphisms transfer through the same
bijections. This proves equivalence of groupoids, without asserting
an unproved equivalence for any broader category of morphisms.
\end{proof}

\begin{counterexample}[Refuted: the membership-free final At\(^*\) list]
\label{SA21-B-Atstar-literal}
Let \(E=\{a,b\}\),
\(\mathscr U=\{\{a\},\{a,b\}\}\), and
\(\mathscr H=\{\varnothing,\mathscr U\}\).
This contains the empty figure and is closed under bounded unions.
The canonical principal subfigures agree wherever they occur.
The only nontrivial inclusion of occurring strata is
\(\{a\}\subset\{a,b\}\), and its canonical singleton figure
refines \(\mathscr U\) in the ambient set-figure atelier.
For the pair \(\mathscr U,\mathscr U\), the premise of At\(^*4\)
fails at the stratum \(\{a\}\), whose canonical singleton is
omitted; for pairs involving the empty figure the conclusion
holds. Finally \(\mathscr U\) is an admissible separating cover,
with residuals \(\{a\},\{b\}\).
Thus the shortened literal conditions hold, but
\(\widehat{\{a\}}=\{\{a\}\}\notin\mathscr H\).
The figure poset has only one nonzero irreducible, so cannot
recover the two prescribed ground objects.
\end{counterexample}

We inspected scan 34 visually. Its At\(^*2\) line asserts only the
equality of canonical principal subfigures; its marginal claim that
At\(^*1\) and At\(^*2\) imply C1--C2 does not supply the missing
membership. The longer formulation at scan 32 explicitly supplies it:
\path{J156-6-04_transcription.tex:189--213,245--249};
\path{J156-6-05_transcription.tex:4--34}.
If the relation in At\(^*3\) is typed only on members of
\(\mathscr H\), the membership is implicit in its typing and the
positive theorem applies. If it is the ambient set-figure relation,
the displayed counterexample applies. We have proved both readings.

\subsection{The Cont and primitive-subdivision systems: exact bridges}

\begin{proposition}[Proved: the final Cont bridge]
\label{SA21-B-Cont-bridge}
The final Cont 1--5 plus Fig 1--4 of 156-4/34--40, with the
geometric-dimension clause retained and with transitivity of internal
refinement added, are equivalent to the prime At 1--7 plus L1,
D\(_0\) and that same dimension clause. The optional finite-principal
clause Cont \(3'\) can be added on both sides. The transitivity
addition cannot be omitted, even if Cont 6 and Cond 7 are imposed.
\end{proposition}

\begin{proof}
The figure clauses are exactly the prime downset incidence system.
Cont 1 identifies incidence with compatible refinement, giving
At 6; refinement of figures is defined by the stratum test At 4.
Cont 4(a) is L1, and Cont 5 is D\(_0\).
It remains to obtain the least carrier. For \(X\ll F\), choose
\(Y\in\widetilde F\) refined by \(X\). Finite height below \(Y\)
gives a minimal \(A\leq Y\) refined by \(X\), so \(X\rho A\).
For any other \(Z\in\widetilde F\) refined by \(X\), choose
similarly \(B\leq Z\) with \(X\rho B\). A place \(p\rho X\)
exists by L1. Internal transitivity gives \(p\in U_A\cap U_B\).
The compatible strata \(A,B\) must coincide by Cont 4(b).
Thus \(A\leq Z\) for every candidate \(Z\), proving At 5.
Figure refinement is a partial order by the carrier proof in
Theorem~\ref{SA21-B-mag-equivalence}, which does not otherwise use
internal disjointness. Conversely, in the prime core all figure
clauses and Cont 1--2 hold. Carrier uniqueness makes open shadows
of distinct compatible strata setwise disjoint, supplying Cont 4(b).
The remaining clauses are exactly the stated added hypotheses.
\end{proof}

\begin{counterexample}[Refuted: Cont 1--5 without internal transitivity]
\label{SA21-B-Cont-counter}
Use the two orders on \(\{p,q,z,u,t\}\) from
Counterexample~\ref{SA21-B-missing-transitivity}, and specify
compatibility independently as equality together with the three
undirected pairs \(pq,qz,zt\). Take all clique downsets as figures.
The only strict incidence is \(z<t\), and it is exactly compatible
refinement. Compatibility is incidence-hereditary. All principal
figures and all their closed subfigures occur; compatible unions
occur. Every refinement has combinatorial dimension at most one.
The open lieu sets listed earlier are nonempty and disjoint on
the three compatible pairs. The only nonzero disjoint elementary
pairs are \(p,q\) and \(q,z\); refining \(z\) can replace it by
\(p\), so disjointness is hereditary. Thus final Cont 1--5 hold.
Whole-shadow disjointness detects exactly those pairs, so Cont 6
holds; open lieux of compatible distinct strata are orthogonal,
so Cond 7 holds. Nevertheless \(p\rho u\rho t\) and
\(p\not\rho t\). The finite certificate checks all figures,
carriers and refinement pairs; it also verifies prime At 1--6.
\end{counterexample}

The early shadow-reflection proposal at 156-4/25 is stronger still:
the two-object local magasin of
Counterexample~\ref{SA21-B-extension-unique} has identical small
multiombres for \(a,x\), although \(x\not\ll a\). Thus even all
the later local and extension-existence axioms do not recover that
early reflection proposal. Conversely strong Mag 1--4 alone do
not impose the nonempty open shadows or finite dimension of Cont;
the preceding countermodels exhibit both failures. The comparison
therefore involves genuine added conditions, not just renumbering.

\begin{proposition}[Proved: the literal C6 uniqueness makes C8\(_0\) redundant]
\label{SA21-B-C6-C8}
Read C6(a) literally: for \(F'\preceq F\) and \(G\leq F\), there is
exactly one \(H\) satisfying \(H\leq F'\) and \(H\preceq G\).
Then
\[
 F\leq H\preceq F\quad\Longrightarrow\quad F=H
\]
for arbitrary figures, and both clauses of C8\(_0\) follow.
With C7\(_0\), the derived relation
\(P\ll F\Longleftrightarrow\exists H\ (P\leq H\preceq F)\)
is a partial order.
\end{proposition}

\begin{proof}
Apply C6(a) to \(H\preceq F\) and the subfigure \(F\leq F\).
Both \(H\) and \(F\) satisfy its displayed conditions when
\(F\leq H\), so uniqueness gives equality. The other C8\(_0\)
configuration is the same inequality with names exchanged.
For transitivity, if \(P\leq A\preceq Q\) and
\(Q\leq B\preceq R\), extend \(A\preceq Q\leq B\)
by C7\(_0\) to \(A\leq B'\preceq B\); then
\(P\leq B'\preceq R\). For antisymmetry suppose
\(P\leq A\preceq Q\) and \(Q\leq B\preceq P\).
Extend \(A\preceq Q\leq B\) to \(A\leq B'\preceq B\).
Then \(P\leq A\leq B'\preceq P\), so \(P=B'\) by the first
assertion, and hence \(A=P\). Thus \(P\preceq Q\).
The inequalities \(Q\leq B\preceq P\preceq Q\) similarly give
\(B=Q\), hence \(Q\preceq P\). Antisymmetry of \(\preceq\)
gives \(P=Q\).
\end{proof}

We inspected the actual scan 156-4/72. The main C6(a) clause says
that there is a unique object satisfying the square; the stratum
formula (b) is a separate marginal addition. The literal quantifier is at
\path{J156-4-09_transcription.tex:161--176}; C7\(_0\) is at
lines 187--191; C8\(_0\) is at
\path{J156-4-11_transcription.tex:49--77}.
We therefore use unrestricted uniqueness exactly as written, rather
than silently replacing it by uniqueness within a smaller class of
candidates. The redundancy conclusion concerns this confirmed text.

\begin{counterexample}[Refuted: even prime C1--C8\(_0\) need not give At 5]
\label{SA21-B-C-carriers}
Let \(M=\{d,a,y,z\}\), with \(d<a,d<y,d<z\).
Every pair is compatible except \(a,y\) and \(a,z\).
The figures, written by their deployments, are
\[
 0,\quad d,\quad da,\quad dy,\quad dz,\quad dyz.
\]
Primitive subdivision is equality together with
\[
 da\preceq dy,\qquad da\preceq dz,\qquad da\preceq dyz.
\]
This is a partial order. All nonempty figures contain \(d\),
so no two are disjoint and C5 holds. C1--C3, including (B.P),
are the explicit clique-downset construction. The incidence
graph of \(M\) is connected, so C4 and C4 bis impose no additional
condition.

For \(da\preceq dy\) or \(da\preceq dz\), restriction sends the
bottom, the subfigure \(d\), and the whole parent respectively to
\(0,d,da\). For \(da\preceq dyz\), it sends
\(0,d,dy,dz,dyz\) respectively to \(0,d,da,da,da\).
These are the unique subfigures that subdivide the prescribed
subfigure; they have exactly the strata refining it, and preserve
all joins. Identity subdivisions restrict identically. This checks
C6(a,b,c) completely. A subdivision \(da\preceq dy\leq dyz\)
extends by \(da\preceq dyz\), and likewise for \(dz\); all other
extension cases are identities. Hence C7\(_0\) holds, as does
C8\(_0\).

Nevertheless \(a\ll dyz\), and the strata refined by \(a\) in
\(dyz\) are \(y,z\), with no least member. Thus the later At 5
fails. Final C8, the \emph{full gluing} axiom for subdivisions,
is a different assertion and fails here: the subdivisions
\(da\preceq dy\) and \(dz\preceq dz\) agree on \(d\) but cannot
glue, because \(a,z\) are incompatible.
The script \texttt{legacy\_subdivision.py} exhaustively checks all
bounded joins, every C6 restriction and every C7\(_0\) extension.
\end{counterexample}

\begin{proposition}[Proved: the modern-to-subdivision bridge]
\label{SA21-B-modern-C}
Suppose a prime atelier has the carrier calculus, L1, L3, C,
At supp, At 14, and gluing for the covers under consideration.
Define subdivision by (B.S). Then it satisfies C1--C3, C5,
literal C6(a,b,c), C7\(_0\), C8\(_0\), and subdivision gluing C8.
The independent incidence-component axioms C4 and C4 bis are not
consequences. In a full canonical figure family C4 implies C4 bis;
for a restricted atlas it need not.
\end{proposition}

\begin{proof}
C1--C3 are assumed; C5 is the definition of subdivision; C7\(_0\)
is At 14; and C8\(_0\) follows from antisymmetry of refinement.
The proof that a subdivision is a maximal refinement uses only
L1, L3 and the relative disjointness calculus: a place in a
missing stratum of a larger compatible refinement is disjoint
from the original one. At supp proves the converse exactly as
in Proposition~\ref{SA21-B-extension-Zorn}.

Let \(F'\preceq F\), \(G\leq F\), and \(R=F'|G\).
If \(R<P\ll G\), then C applied to the compatible parents
\(G\leq F\) makes \(P\) compatible with \(F'\), because their
restrictions are \(P\) and \(R\). Their union would be a strictly
larger refinement of \(F\), contradicting maximality of \(F'\).
Thus \(R\) is a maximal refinement of \(G\), hence a subdivision.
If \(H\leq F'\) also subdivides \(G\), its strata refine \(G\),
so \(H\leq R\). Maximality of \(H\) forces \(H=R\).
This is literal C6(a); (b) and (c) are the carrier inverse-image
formula and its union preservation.

Gluing a matching family of subdivisions gives a refinement \(P\)
of the covered figure. Any larger refinement restricts to a
larger refinement of each subdivision, hence to the same one by
maximality. The cover's carrier formula then forces equality
globally. Thus \(P\) is maximal and is a subdivision.
For the last assertion, C4 makes the union of figures chosen in
distinct incidence components a clique downset. The full figure
family admits it. A restricted atlas may omit it, as the next
example shows.
\end{proof}

\begin{counterexample}[Refuted: C4 does not force C4 bis in restricted atlases]
\label{SA21-B-C4bis}
Take an infinite set of independent point objects, both orders
equal to identity, and internal disjointness equal to inequality.
Admit only finite figures. All the strong local axioms, finite
gluing and At Fil hold; the latter holds because its parent has
only finitely many strata, so local stationarity is joint
stationarity. Every incidence component is one point and distinct
components are compatible, so C4 holds. The family containing
one point figure from every component has no figure union, so
C4 bis fails. This failure concerns an unbounded family and is
consistent with At Fil.
\end{counterexample}

\begin{counterexample}[Refuted: cosupport preservation does not characterize primitive subdivision]
\label{SA21-B-primitive-not-support}
There is a second reason not to identify primitive subdivision
with (B.S) without an extra axiom. For the two-element incidence
chain \(a<x\), with all downsets as figures and primitive
subdivision equal to identity, C5, C6, C7\(_0\) and C8 hold
tautologically. Derived refinement is incidence. Both nonempty
figures have cosupport \(\{0\}\), so (B.S) would additionally
declare \(a\preceq x\). C5 says that subdivisions preserve
cosupport; it does not say every same-support refinement is one.
The complete two-stratum structure, including every C6 restriction,
C7 extension and C8 gluing assignment, is checked in
\texttt{legacy\_subdivision.py}.
\end{counterexample}

\subsection{The 156-9 spatial presentation and the comparison ledger}

The typed spatial reconstruction is proved in
Theorem~\ref{SA21-A23}. Its use in the axiom comparison is precise.
At spat 1 gives incidence and symmetric irreflexive internal
disjointness. At spat 2 is the comparable-pair implication. At
spat 3 places each deployment, injectively, in the space of
\emph{ordered spatial figures}. On scan 48 that target has two
explicit conditions: every label is nonzero, and distinct labels
are orthogonal. For principal labels, orthogonality is exactly
internal disjointness of their generators. Consequently the
target membership supplies (B.M4); it cannot be replaced merely
by injectivity into ordered families of disjoint sets.
The source itself requests the additional At 4 in the margins
of scans 47 and 49:
\path{J156-9-06_transcription.tex}, scans 47--49, and
\path{J156-9-07_transcription.tex:11--39}, scan 50.
The five-cycle example in Theorem~\ref{SA21-A23} verifies the
failure of that weaker untyped reading. The change from
``magasin'' to ``atelier'' on scan 41 changes terminology,
not these relations.

\begin{counterexample}[Refuted: spatiality and divisibility are automatic]
\label{SA21-B-spatial-separations}
Strong Mag 1--4, even together with Mog L2 and all the extension
existence axioms, do not force At spat 3. In the two-object model of
Counterexample~\ref{SA21-B-extension-unique}, internal orthogonality
is empty, so the orthoclosure of either singleton is \(M\). The two
elementary deployments have the same single label \(M\) and the
same one-point order. They therefore fail the required injectivity.

Conversely the typed spatial axioms do not force divisibility.
Take incidence \(a<t\), \(\rho=\mathrm{id}\), and
\(\delta\) equal to inequality on \(\{a,t\}\). Orthoclosures are
singletons. The two deployments are the labelled singleton
\(\{\{a\}\}\) and the labelled chain
\(\{a\}<\{t\}\); labels within a deployment are nonzero and
orthogonal, and the deployments are distinct. A top-preserving
carrier by inclusion of these singleton labels exists exactly for
an identical source and target, recovering \(\rho=\mathrm{id}\).
Thus the fully typed spatial axioms hold. The only lieu is \(a\),
but \(U_t=\varnothing\), so divisibility fails. The finite
certificate checks the orthoclosures and both deployment assertions.
\end{counterexample}

\begin{longtable}{p{.42\linewidth}p{.47\linewidth}}
\hline
Proposed identification or converse & Exact outcome and witness\\ \hline
Literal irreducible C1--C3 = prime downset system & False;
\(M_3\), Counterexample~\ref{SA21-B-M3}. Exact repair (B.P),
Theorem~\ref{SA21-B-prime-representation}.\\
Final Cont = later core & Add internal transitivity;
Proposition~\ref{SA21-B-Cont-bridge}. The five-object
Counterexample~\ref{SA21-B-Cont-counter} shows necessity.\\
C1--C8\(_0\) = later carrier core & False even after primality;
Counterexample~\ref{SA21-B-C-carriers}. Final C8 gluing is a
different extra axiom.\\
Strong Mag = corrected strong primed list & True by
Theorem~\ref{SA21-B-mag-equivalence}; comparable-only list
is weaker by Counterexample~\ref{SA21-B-weak-primed}.\\
Strong Mag = local lieu presentation & Add Mog L2;
Theorem~\ref{SA21-B-local-equivalence}. Divisibility alone
does not suffice, Counterexample~\ref{SA21-B-L6-countermodels}(1).\\
Strong Mag implies all optional criteria & False;
Counterexamples~\ref{SA21-B-L6-countermodels}--\ref{SA21-B-fil}.
The valid relative and atlas-closure clauses are stated in
Remark~\ref{SA21-B-L6-status}.\\
C implies D implies D\(_0\) & True; neither converse holds.
Use Counterexamples~\ref{SA21-B-Sigma-counter} and
\ref{SA21-B-gluing-not-L3}, respectively.\\
Finite gluing implies C or L3 & False; use
Counterexamples~\ref{SA21-B-Sigma-counter} and
\ref{SA21-B-gluing-not-L3}, respectively.\\
Finite gluing implies arbitrary gluing & False for restricted
atlases, Counterexample~\ref{SA21-B-fil}; At Fil supplies
the exact locally stationary union step.\\
At \(9_\Sigma\)+At 10 implies At 12 & False even with At 11
and D. Exact replacement is (B.Gate).\\
At 14 bis's diagram implies unique object & False in two
objects; only the abstract deployment is unique.\\
The final At\(^*\) list reconstructs its ground objects & True
with canonical-subfigure membership and a prime source core;
the membership-free reading has a two-point counterexample.\\
Spatial At spat 1--3 = ordinary set-label injectivity & False;
the orthogonal spatial target is indispensable, as in
Theorem~\ref{SA21-A23}.\\
Strong Mag implies spatiality, or spatiality implies divisibility &
Both false; the two-object models of
Counterexample~\ref{SA21-B-spatial-separations} separate them.\\ \hline
\end{longtable}

\section{B.2. What replaces the false global connectedness axiom}
\label{SA21-B-connectedness}

The global incidence graph of \(M\), the incidence graph of a
\emph{single figure}, and the global refinement graph are three
different objects. We give their exact relationships. In
particular, replacing the first by either of the latter two is
a change of statement that must be made explicit.

\begin{theorem}[Proved: all incidence components of \(M(\operatorname{Fig}(L))\)]
\label{SA21-B-incidence-components}
Let \(L\) be an arbitrary set. A multistrate in
\(\operatorname{Fig}(L)\) is an elementary admissible family of
subsets of \(L\), ordered by closed-subfamily incidence. If
\(|L|\geq3\), the connected components of its comparability
graph are exactly:
\begin{enumerate}
\item the isolated multistrate \(\{L\}\);
\item for each \(p\in L\), the two-object component
\[
 \{L\setminus\{p\}\}
 \ <\ \{L\setminus\{p\},L\};
\]
\item one component containing every remaining multistrate.
\end{enumerate}
The last component contains every point multistrate.
For \(|L|=2\) there are just the isolated object and the two
two-object components. For \(|L|=1\) there is one object and one
component; for \(L=\varnothing\), there are no multistrates.
Thus a finite \(n\)-point set has \(n+2\) components for
\(n\geq3\), three for \(n=2\), one for \(n=1\), and none for
\(n=0\).
\end{theorem}

\begin{proof}
The multistrate \(\{L\}\) has no proper elementary subfigure.
It has no proper elementary superfigure: its sole stratum is
already the largest possible support, and closed-subfamily
incidence forbids adding lower strata beneath it.

Put \(A=L\setminus\{p\}\). If an elementary superfigure
contains \(\{A\}\) as a closed subfigure, no new stratum can lie
properly in \(A\). Its top is either \(A\) or \(L\).
In the latter case any other new proper stratum would contain
\(p\), and together with \(A\) would exhaust \(L\), leaving
the top with empty residual. Thus the only possible strict
superfigure is \(\{A,L\}\), whose only proper elementary
subfigure is \(\{A\}\). This proves the claimed exceptional
components.

Suppose an elementary figure \(\Phi\), of top support \(A\),
has two distinct unused points \(y,z\in L\setminus A\).
Then
\[
 \Psi=\Phi\cup\bigl\{\{y\},\ A\cup\{y,z\}\bigr\}
\]
is admissible and elementary: the old residuals are unchanged,
the new point residual is \(\{y\}\), and the new top residual
is \(\{z\}\). Both \(\Phi\) and the point figure \(\{\{y\}\}\)
are elementary subfigures of \(\Psi\). Thus \(\Phi\) lies in
a component containing a point. Any two distinct point figures
are connected using a third point \(w\), through
\(\{\{u\},\{v\},\{u,v,w\}\}\).

Finally consider any nonexceptional elementary \(\Phi\).
If one of its strata \(B\) has at least two unused points,
its principal subfigure \(\Phi_B\leq\Phi\) has just been linked
to a point. Otherwise every stratum is \(L\) or a co-singleton.
There can be at most one proper co-singleton below the top:
two distinct ones cover \(L\), again making its residual empty.
The only resulting elementary families are the exceptions
already listed. This contradiction proves that every
nonexceptional figure lies in the single remaining component.
The argument works without any finiteness assumption on \(L\)
or on the admissible families. The cases of size at most two
follow by direct enumeration of their nonempty subsets.
\end{proof}

\begin{corollary}[Refuted: C4 in the full set-theoretic model]
\label{SA21-B-C4-false}
For every set \(L\) with at least two elements, C4 fails in
\(\operatorname{Fig}(L)\). The isolated object \(\{L\}\)
and any point figure lie in different incidence components
and are incompatible. Their underlying supports meet, while
their common incidence subfigure is empty.
\end{corollary}

This independently rechecks input result SA21-K5-01 and corrects
the claimed connecting constructions of 156-5/13--15,
\path{J156-5-02_transcription.tex:70--115}: their
proposed top set is exhausted by the lower sets and therefore
has empty residual. The script \texttt{connectedness.py}
enumerates every admissible elementary family for \(0\leq n\leq4\),
forms every principal-subfigure edge, and verifies the complete
component partition, not only the number of components:
\[
\begin{array}{c|r|r|l}
 |L|&|M|&\#\text{components}&\text{component sizes}\\ \hline
 0&0&0&\text{none}\\
 1&1&1&1\\
 2&5&3&1,2,2\\
 3&28&5&1,2,2,2,21\\
 4&231&6&1,2,2,2,2,222.
\end{array}
\]

\begin{theorem}[Proved: the correct connectedness of each figure]
\label{SA21-B-figure-components}
Let \(\Phi\) be any admissible family on \(L\), and let
\(\pi:|\Phi|\to\Phi\) send a point to its least stratum.
Give \(|\Phi|\) the Alexandrov topology whose closed sets are
\(\pi^{-1}(D)\) for downsets \(D\subseteq\Phi\).
Its connected components are exactly the inverse images of
the comparability components of the stratum poset \(\Phi\).
No finiteness hypothesis is needed.

For a finite frontier stratification by definably connected
cells, with each closed stratum the closure of its cell, these
same components are exactly the definable connected components
of its actual support. In particular this applies to the
o-minimal cell atelier of Part B.4.
\end{theorem}

\begin{proof}
Every residual fiber \(\pi^{-1}(A)\) is nonempty and has the
indiscrete subspace topology. A comparability component of
\(\Phi\) is both a downset and an upset, so its inverse image
is clopen. Any clopen subset must be constant on each fiber,
and cannot separate comparable fibers. Finite comparability
zigzags therefore cannot be separated. The stated clopen
pieces are connected and maximal.

In the definable model, a poset component is a finite union
of closed strata; its complement in the support is another
such union. Hence it is definably clopen. Conversely a
definably clopen subset cannot split a definably connected
cell. If it contains a cell, closedness makes it contain
that cell's entire closure, hence every incident lower cell.
Applying the same reasoning to its complement forces it to
be a union of comparability components. This proves equality
of the two component decompositions.
\end{proof}

The finite-poset part is the standard finite-space connectedness
criterion, which we checked in
\citet[Proposition 1.2.4, p.~5]{Barmak2011}; our proof also gives
the possibly infinite admissible-family version and explicitly
handles nonsingleton residual fibers. Over a non-Archimedean real
closed field, the geometric assertion concerns \emph{definable}
connectedness, not arbitrary topological connectedness.

\begin{theorem}[Proved: the global refinement replacement for C4]
\label{SA21-B-refinement-components}
Let \(\mathscr A\) be a set-theoretic atelier on \(L\), with
inherited set-figure refinement, containing every point figure
and closed under disjoint unions. Make a hypergraph on \(L\)
whose hyperedges are the supports of the multistrates.
The comparability components of \((M,\ll)\) are exactly the
components of that support hypergraph, with a multistrate
assigned to the component containing its support.
Multistrates in different refinement components are compatible.

For the full \(\operatorname{Fig}(L)\), the refinement graph
is connected whenever \(L\ne\varnothing\). For the closed
o-minimal cell atelier of Part B.4, its refinement components
are exactly the definable connected components of the ambient
definable set \(X\).
\end{theorem}

\begin{proof}
Every \(p\in|A|\) supplies the refinement \(\{p\}\ll A\).
Thus all points of one hyperedge are joined through its
multistrate. Conversely \(A\ll B\) implies
\(|A|\subseteq|B|\), so a refinement edge never crosses a
hypergraph component. These two observations prove equality
of the component relations. Distinct components have disjoint
supports, and closure under disjoint unions gives compatibility.
For the full set-theoretic atelier, any two points lie in the
one-stratum multistrate \(\{\{p,q\}\}\), so there is one component.

In the o-minimal model, every multistrate support is the closure
of a definably connected cell, hence is definably connected.
It lies in one definable component of \(X\), so refinement
components cannot be larger. If \(p,q\) lie in one definable
component, \citet[Theorem 3.9, p.~28]{Coste1999} supplies a
definable continuous path \(\gamma:[0,1]\to X\) between them.
Its image \(P\) is closed and bounded: coordinate functions are
bounded, and for a putative missing limit point \(v\), the positive
function \(t\mapsto\|\gamma(t)-v\|\) attains a positive minimum,
contradicting that \(v\) is a limit point. These uses of definable
compactness follow from \citet[Theorem 3.4, pp.~26--27]{Coste1999}.
Triangulate \(P\) compatibly with \(\{p\},\{q\}\), using
\citet[Theorem 4.4, p.~42]{Coste1999}. The resulting finite closed
cell figure has connected support, so its incidence graph is
connected by Theorem~\ref{SA21-B-figure-components}. Incidence
implies refinement, and its endpoint point figures therefore
lie in one refinement component. This proves the converse.
No injective-arc theorem is needed.
\end{proof}

The original C4 still fails in the geometric model. For example,
with \(X=R\), the single cell \(R\) has an isolated multistrate
\(\{R\}\) in the global incidence graph. It is incompatible
with the point figure \(\{\{0\}\}\), which lies in another
incidence component. Thus the two positive theorems above are
the corrected statements, and do not reassert the false global
incidence claim under a different name.
